% =====================================================================
% Chủ đề 5 — Tư duy logic và Câu đố — phần b: bài 14–30 (bài mới, mức cơ bản)
% =====================================================================
% quả cân có ghi khối lượng (Bài 22)
\providecommand{\lgaWeight}[1]{\tikz[baseline=(lgaw.south)]{\node[trapezium,trapezium angle=72,draw=Slate,fill=Slate!25,line width=0.6pt,
  rounded corners=1.5pt,inner xsep=1pt,inner ysep=2.5pt,font=\scriptsize\bfseries] (lgaw) {#1};}}

% ---------------------------------------------------------------- Bài 14
\needspace{12\baselineskip}\begin{probbox}{14}{\relax}
Trên một hòn đảo, mỗi người hoặc luôn nói thật, hoặc luôn nói dối. Em gặp hai người là Alice và Ben. Alice nói:
``Cả hai chúng tôi đều là người nói dối.'' Khẳng định nào sau đây đúng?
\choicesL{Cả hai đều nói thật.}{Cả hai đều nói dối.}{Alice nói thật, Ben nói dối.}{Alice nói dối, Ben nói thật.}{Không xác định được.}
\end{probbox}
\probsol{
\textbf{Phân tích đề.} Chỉ có một câu nói, của Alice. Ta xét hai trường hợp: Alice nói thật hoặc Alice nói dối, rồi xem trường hợp nào
không mâu thuẫn.
\begin{recallbox}Người nói thật chỉ nói câu đúng; người nói dối chỉ nói câu sai. Phủ định của ``cả hai đều nói dối'' là
``ít nhất một người nói thật''.\end{recallbox}
\step{1}{Trường hợp Alice nói thật}
Khi đó câu của Alice đúng: cả hai đều nói dối -- tức là Alice cũng nói dối. Nhưng ta đang giả sử Alice nói thật. \textbf{Mâu thuẫn}, loại.
\step{2}{Trường hợp Alice nói dối}
Khi đó câu của Alice sai: \emph{không phải} cả hai đều nói dối, tức là có ít nhất một người nói thật. Alice nói dối, nên người nói thật
phải là \textbf{Ben}. Không có mâu thuẫn.
\step{3}{Kết luận}
Alice nói dối, Ben nói thật.
\begin{proofbox}Đây là một dạng của câu ``Tôi là người nói dối'': một người nói thật không bao giờ tự nhận mình nói dối. Vì vậy ngay khi
nghe Alice nói ``chúng tôi đều nói dối'' (có cả Alice trong đó), ta biết Alice phải là người nói dối.\end{proofbox}
\begin{warnbox}Nhiều bạn chọn (B) vì nghĩ ``Alice nói dối nên điều ngược lại là cả hai đều nói thật''. Sai! Điều ngược lại của ``cả hai đều
nói dối'' chỉ là ``ít nhất một người nói thật''.\end{warnbox}
\verify{Alice dối, Ben thật: câu ``cả hai đều nói dối'' sai (vì Ben nói thật) -- đúng là một câu nói dối, phù hợp với Alice.}
Vậy chọn \textbf{(D)}.
\ans{(D) Alice nói dối, Ben nói thật}
}

% ---------------------------------------------------------------- Bài 15
\needspace{12\baselineskip}\begin{probbox}{15}{\relax}
Một chiếc bình hoa bị vỡ. Có đúng một trong bốn bạn David, Helen, Leo, Mia đã làm vỡ bình. Khi được hỏi, các bạn trả lời:
\begin{itemize}[leftmargin=1.6em,itemsep=0pt,topsep=2pt]
\item David: ``Helen làm vỡ.'' \qquad Helen: ``Leo làm vỡ.''
\item Leo: ``Helen nói dối.'' \qquad Mia: ``Không phải tôi làm vỡ.''
\end{itemize}
Biết rằng chỉ có \textbf{đúng một} bạn nói thật. Ai đã làm vỡ bình hoa?
\choicesII{David}{Helen}{Leo}{Mia}{Không xác định được}
\end{probbox}
\probsol{
\textbf{Phân tích đề.} Bốn câu nói, đúng một câu đúng, ba câu sai. Cần tìm người làm vỡ bình.
\begin{recallbox}Hai câu mâu thuẫn (không thể cùng đúng, không thể cùng sai) thì đúng một trong hai câu đúng.\end{recallbox}
\step{1}{Tìm cặp câu mâu thuẫn}
Leo nói ``Helen nói dối''. Nếu câu của Helen đúng thì câu của Leo sai; nếu câu của Helen sai thì câu của Leo đúng. Vậy trong hai câu của
Helen và Leo có \textbf{đúng một} câu đúng.
\step{2}{Suy ra David và Mia đều nói dối}
Cả nhóm chỉ có một câu đúng, và câu đúng đó nằm trong cặp Helen -- Leo. Do đó câu của David và câu của Mia đều \textbf{sai}.
\step{3}{Dùng câu sai của Mia}
Mia nói ``Không phải tôi làm vỡ'' là câu sai, nên \textbf{Mia đã làm vỡ bình}.
\step{4}{Thử lại toàn bộ}
\begin{center}
\begin{tabular}{|l|l|c|}
\hline
\rowcolor{tmBlue!10} Bạn & Câu nói & Đúng / sai (khi Mia làm vỡ) \\ \hline
David & ``Helen làm vỡ'' & sai \\ \hline
Helen & ``Leo làm vỡ'' & sai \\ \hline
Leo & ``Helen nói dối'' & \textbf{đúng} (Helen đúng là nói sai) \\ \hline
Mia & ``Không phải tôi'' & sai \\ \hline
\end{tabular}
\end{center}
Đúng một câu đúng (của Leo). Khớp với đề.
\begin{warnbox}Đừng thử ngẫu nhiên từng người rồi dừng lại khi thấy ``có vẻ ổn''. Cách chắc chắn là thử đủ $4$ khả năng: nếu David làm vỡ
thì câu của Leo và Mia đều đúng ($2$ câu đúng); nếu Helen làm vỡ thì cả ba câu của David, Leo và Mia đều đúng; nếu Leo làm vỡ thì câu của Helen và Mia đúng.
Chỉ ``Mia làm vỡ'' cho đúng một câu đúng.\end{warnbox}
Vậy chọn \textbf{(D)}.
\ans{(D) Mia}
}

% ---------------------------------------------------------------- Bài 16
\needspace{12\baselineskip}\begin{probbox}{16}{\relax}
Bạn Kate dùng mật mã Caesar với bước dịch $5$ để mã hoá một lời chúc tiếng Anh (mỗi chữ cái được thay bằng chữ cái đứng sau nó
$5$ vị trí trên vòng chữ cái A--Z, sau Z lại đến A). Kate nhận được: \textbf{MFAJ KZS}. Lời chúc ban đầu là
\choicesII{HAVE FUN}{RKFO PEX}{IBWF GVO}{HAVE FUM}{GZUD ETM}
\end{probbox}
\probsol{
\textbf{Phân tích đề.} Ta có bản mã và biết bước dịch $5$. Muốn tìm bản gốc, ta phải \emph{giải mã}: dịch \emph{ngược} $5$ vị trí.
\begin{recallbox}Giải mã: lấy số thứ tự của chữ (A $= 0$, \ldots, Z $= 25$) trừ đi $5$; nếu được số âm thì cộng thêm $26$.\end{recallbox}
\step{1}{Giải mã từng chữ}
\begin{center}
\setlength{\tabcolsep}{4pt}
\begin{tabular}{|c|c|c|c|c|c|c|c|}
\hline
\rowcolor{tmBlue!10} Chữ mã & M & F & A & J & K & Z & S \\ \hline
Số & $12$ & $5$ & $0$ & $9$ & $10$ & $25$ & $18$ \\ \hline
Trừ $5$ & $7$ & $0$ & $-5 \to 21$ & $4$ & $5$ & $20$ & $13$ \\ \hline
Chữ gốc & \textbf{H} & \textbf{A} & \textbf{V} & \textbf{E} & \textbf{F} & \textbf{U} & \textbf{N} \\ \hline
\end{tabular}
\end{center}
Chữ A phải ``quay vòng'': A $\to$ Z $\to$ Y $\to$ X $\to$ W $\to$ V (lùi $5$ bước), đúng bằng $-5 + 26 = 21 =$ V.
\step{2}{Đọc kết quả}
Lời chúc là HAVE FUN (tiếng Anh: ``chúc vui vẻ'').
\begin{warnbox}(B) RKFO PEX là kết quả khi dịch \emph{xuôi} thêm $5$ bước -- tức là mã hoá thêm một lần, chứ không phải giải mã. (C) và (E)
là dịch ngược $4$ bước và $6$ bước (đếm sai). (D) sai ở chữ cuối: S lùi $5$ bước là N, không phải M.\end{warnbox}
\verify{Mã hoá HAVE FUN với bước $5$: H $\to$ M, A $\to$ F, V$(21) + 5 = 26 - 26 = 0 =$ A, E $\to$ J, F $\to$ K, U $\to$ Z, N $\to$ S.
Được MFAJ KZS.}
Vậy chọn \textbf{(A)}.
\ans{(A) HAVE FUN}
}

% ---------------------------------------------------------------- Bài 17
\needspace{12\baselineskip}\begin{probbox}{17}{\relax}
Bạn Nora mã hoá chữ M bằng mật mã Caesar ba lần liên tiếp: lần thứ nhất với bước dịch $9$, lần thứ hai với bước dịch $11$ (dịch tiếp
từ kết quả lần một), lần thứ ba lại với bước dịch $9$. Sau ba lần, chữ M trở thành chữ nào?
\choices{V}{G}{J}{P}{O}
\end{probbox}
\probsol{
\textbf{Phân tích đề.} Ba lần dịch theo cùng chiều. Ta có thể dịch từng lần, hoặc gộp lại thành một lần dịch.
\begin{recallbox}Dịch $a$ bước rồi dịch $b$ bước giống như dịch $a + b$ bước. Dịch $26$ bước thì về chỗ cũ, nên chỉ cần số dư khi chia
số bước cho $26$.\end{recallbox}
\step{1}{Gộp ba lần dịch}
Tổng số bước: $9 + 11 + 9 = 29$. Vì $29 = 26 + 3$, dịch $29$ bước giống hệt dịch $3$ bước.
\step{2}{Dịch chữ M đi $3$ bước}
M $= 12$; $12 + 3 = 15 =$ P. (M $\to$ N $\to$ O $\to$ P.)
\step{3}{Kiểm tra bằng cách dịch từng lần}
\begin{center}
\begin{tabular}{|c|c|c|c|}
\hline
\rowcolor{tmBlue!10} Ban đầu & Sau lần 1 ($+9$) & Sau lần 2 ($+11$) & Sau lần 3 ($+9$) \\ \hline
M $(12)$ & $21 =$ V & $32 - 26 = 6 =$ G & $15 =$ P \\ \hline
\end{tabular}
\end{center}
\begin{warnbox}Các phương án (A) V và (B) G là kết quả \emph{giữa chừng} (sau lần $1$ và lần $2$). Đọc kĩ: đề hỏi kết quả sau \emph{cả ba} lần.
(C) J là dịch \emph{ngược} $3$ bước.\end{warnbox}
\verify{Từ P lùi $29$ bước, tức lùi $3$ bước (vì lùi $26$ bước là về chỗ cũ): P $\to$ O $\to$ N $\to$ M. Đúng.}
Vậy chọn \textbf{(D)}.
\ans{(D) P}
}

% ---------------------------------------------------------------- Bài 18
\needspace{12\baselineskip}\begin{probbox}{18}{\relax}
Năm bạn Nina, Olivia, Peter, Quinn, Sam có chiều cao khác nhau. Nina cao hơn Olivia. Peter thấp hơn Quinn nhưng cao hơn Nina. Sam thấp
hơn Olivia. Khi xếp năm bạn theo chiều cao, bạn nào đứng \emph{chính giữa}?
\choices{Nina}{Olivia}{Peter}{Quinn}{Sam}
\end{probbox}
\probsol{
\textbf{Phân tích đề.} Có năm bạn; ``chính giữa'' là bạn cao thứ ba (có hai bạn cao hơn và hai bạn thấp hơn).
\step{1}{Kí hiệu hoá (``$>$'' là ``cao hơn'')}
$N > O$; \quad $Q > P > N$; \quad $O > S$.
\step{2}{Nối thành một dãy}
$Q > P > N$ và $N > O > S$ cho: \quad $Q > P > N > O > S$.
\begin{figC}
  \draw[->,Navy,line width=0.8pt] (0,-0.3)--(0,2.8) node[above,font=\scriptsize,text=Slate]{cao hơn};
  \foreach \y/\n in {2.4/Quinn,1.8/Peter,1.2/Nina,0.6/Olivia,0/Sam}{
    \fill[tmBlue] (0,\y) circle (2.2pt); \node[right,font=\small] at (0.15,\y){\n};}
  \node[anchor=west,font=\small\bfseries,text=AnsRed] at (1.45,1.2){$\leftarrow$ đứng giữa (thứ $3$)};
\end{figC}
\step{3}{Đọc kết quả}
Thứ tự từ cao đến thấp: Quinn, Peter, \textbf{Nina}, Olivia, Sam. Bạn đứng giữa là Nina.
\begin{warnbox}Câu ``Peter thấp hơn Quinn nhưng cao hơn Nina'' có \emph{hai} thông tin. Nhiều bạn chỉ ghi $Q > P$ mà quên $P > N$, khi đó
không xếp được dãy.\end{warnbox}
\verify{Nina (thứ $3$) cao hơn Olivia (thứ $4$) \checkmark; Peter (thứ $2$) thấp hơn Quinn, cao hơn Nina \checkmark; Sam (thứ $5$) thấp hơn Olivia \checkmark.}
Vậy chọn \textbf{(A)}.
\ans{(A) Nina}
}

% ---------------------------------------------------------------- Bài 19
\needspace{12\baselineskip}\begin{probbox}{19}{\relax}
Trong một hàng người xếp hàng mua vé, bạn Henry đứng thứ $7$ tính từ đầu hàng và đứng thứ $12$ tính từ cuối hàng. Hàng đó có bao nhiêu người?
\choices{$17$}{$18$}{$19$}{$20$}{$5$}
\end{probbox}
\probsol{
\textbf{Phân tích đề.} Biết vị trí của Henry khi đếm từ hai đầu. Cần tìm tổng số người trong hàng.
\step{1}{Vẽ sơ đồ}
\begin{figC}
  \foreach \i in {1,...,18}{\fill[tmBlue!60] (\i*0.5,0) circle (0.15);}
  \fill[AnsRed] (3.5,0) circle (0.17);
  \node[font=\scriptsize,text=AnsRed] at (3.5,0.4){Henry};
  \draw[decorate,decoration={brace,amplitude=4pt,mirror},tmGreen!60!black,line width=0.8pt] (0.35,-0.3)--(3.65,-0.3) node[midway,below=4pt,font=\scriptsize]{$7$ người (tính cả Henry)};
  \draw[decorate,decoration={brace,amplitude=4pt},tmOrange!80!black,line width=0.8pt] (3.35,0.75)--(9.15,0.75) node[midway,above=4pt,font=\scriptsize]{$12$ người (tính cả Henry)};
  \node[font=\scriptsize,text=Slate] at (0.5,0.4){đầu}; \node[font=\scriptsize,text=Slate] at (9,0.4){cuối};
\end{figC}
\step{2}{Đếm}
Đếm từ đầu đến Henry: $7$ người. Đếm từ Henry đến cuối: $12$ người. Cộng lại $7 + 12 = 19$, nhưng Henry được đếm \textbf{hai lần}.
Vậy số người trong hàng là $19 - 1 = 18$.
\begin{recallbox}Đứng thứ $k$ từ đầu và thứ $m$ từ cuối thì hàng có $k + m - 1$ người.\end{recallbox}
\step{3}{Cách khác}
Phía trước Henry có $7 - 1 = 6$ người, phía sau Henry có $12 - 1 = 11$ người. Cả hàng: $6 + 1 + 11 = 18$ người.
\begin{warnbox}Phương án (C) $19 = 7 + 12$ là lỗi đếm Henry hai lần; phương án (E) $5 = 12 - 7$ là lấy hiệu -- không có ý nghĩa gì ở đây.\end{warnbox}
\verify{Hàng $18$ người, Henry đứng thứ $7$ từ đầu thì tính từ cuối là thứ $18 - 7 + 1 = 12$. Đúng.}
Vậy chọn \textbf{(B)}.
\ans{(B) $18$}
}

% ---------------------------------------------------------------- Bài 20
\needspace{12\baselineskip}\begin{probbox}{20}{\relax}
\begin{minipage}[t]{0.63\linewidth}
Điền các số $1, 2, 3, \ldots, 9$ (mỗi số một lần) vào bảng bên cạnh để được một ma phương (tổng các số trên mỗi hàng, mỗi cột, mỗi đường
chéo đều bằng nhau). Ba số đã được điền sẵn. Hãy hoàn thành bảng. Số ở góc dưới bên trái là bao nhiêu?
\end{minipage}\hfill
\begin{minipage}[t]{0.33\linewidth}\vspace{-4pt}\centering
\lgaMagic{$2$}{$7$}{}{}{}{$1$}{}{}{}
\end{minipage}
\end{probbox}
\probsol{
\textbf{Phân tích đề.} Ma phương $3\times 3$ với các số $1$ đến $9$; đã biết ô góc trên trái ($2$), ô giữa hàng trên ($7$), ô bên phải
hàng giữa ($1$). Cần điền đủ $6$ ô còn lại.
\step{1}{Tổng ma phương và số ở giữa}
\begin{recallbox}$S = (\text{tổng } 9 \text{ số}) : 3$; số ở giữa $= S : 3$.\end{recallbox}
$1 + 2 + \cdots + 9 = 45$, nên $S = 45 : 3 = 15$; số ở giữa $= 15 : 3 = 5$.
\step{2}{Điền lần lượt những đường đã có hai số}
\begin{itemize}[leftmargin=1.4em,itemsep=1pt,topsep=1pt]
\item Hàng trên: $2 + 7 + \square = 15$ $\Rightarrow$ góc trên phải $= 6$.
\item Hàng giữa: $\square + 5 + 1 = 15$ $\Rightarrow$ ô trái hàng giữa $= 9$.
\item Đường chéo $2$ -- $5$ -- ?: góc dưới phải $= 15 - 2 - 5 = 8$.
\item Đường chéo $6$ -- $5$ -- ?: góc dưới trái $= 15 - 6 - 5 = \mathbf{4}$.
\item Cột giữa: $7 + 5 + \square = 15$ $\Rightarrow$ ô giữa hàng dưới $= 3$.
\end{itemize}
\begin{center}\lgaMagic{$2$}{$7$}{\textcolor{tmGreen!60!black}{$6$}}{\textcolor{tmGreen!60!black}{$9$}}{\textcolor{tmGreen!60!black}{$5$}}{$1$}{\textcolor{AnsRed}{$\mathbf{4}$}}{\textcolor{tmGreen!60!black}{$3$}}{\textcolor{tmGreen!60!black}{$8$}}\end{center}
\verify{Hàng: $2 + 7 + 6 = 15$, $9 + 5 + 1 = 15$, $4 + 3 + 8 = 15$. Cột: $2 + 9 + 4 = 15$, $7 + 5 + 3 = 15$, $6 + 1 + 8 = 15$.
Chéo: $2 + 5 + 8 = 15$, $6 + 5 + 4 = 15$. Các số $1$ đến $9$ dùng đúng một lần.}
\begin{warnbox}Phải thử lại \emph{tất cả} tám đường (ba hàng, ba cột, hai chéo) và kiểm tra không có số nào bị lặp. Một bảng chỉ đúng
hàng và cột nhưng sai đường chéo thì chưa phải ma phương.\end{warnbox}
\ans{Bảng: $2, 7, 6$ / $9, 5, 1$ / $4, 3, 8$. Số ở góc dưới bên trái là $4$.}
}

% ---------------------------------------------------------------- Bài 21
\needspace{12\baselineskip}\begin{probbox}{21}{\relax}
\begin{minipage}[t]{0.63\linewidth}
Các số $3, 4, 5, \ldots, 11$ (chín số tự nhiên liên tiếp) được điền vào bảng bên cạnh, mỗi số một lần, để được một ma phương.
Số $x$ ở góc dưới bên phải là
\choicesII{$6$}{$8$}{$10$}{$3$}{$5$}
\end{minipage}\hfill
\begin{minipage}[t]{0.33\linewidth}\vspace{-4pt}\centering
\lgaMagic{$4$}{$9$}{}{$11$}{}{}{}{}{$x$}
\end{minipage}
\end{probbox}
\probsol{
\textbf{Phân tích đề.} Chín số từ $3$ đến $11$. Đã biết ba ô; ô $x$ nằm trên đường chéo chính cùng với ô $4$ và ô giữa.
\step{1}{Tìm tổng ma phương}
Tổng chín số: $3 + 4 + \cdots + 11$. Ghép cặp đầu -- cuối: $(3 + 11) + (4 + 10) + (5 + 9) + (6 + 8) + 7 = 4\times 14 + 7 = 63$.
Vậy $S = 63 : 3 = 21$.
\step{2}{Số ở giữa}
\begin{recallbox}Trong ma phương $3\times 3$, số ở giữa $= S : 3$.\end{recallbox}
Số ở giữa $= 21 : 3 = 7$ (cũng chính là số đứng giữa dãy $3, 4, \ldots, 11$).
\step{3}{Tìm $x$ trên đường chéo chính}
$4 + 7 + x = 21$ $\Rightarrow$ $x = 21 - 11 = 10$.
\step{4}{Hoàn thành bảng để thử lại}
Hàng trên: $21 - 4 - 9 = 8$. Hàng giữa: $21 - 11 - 7 = 3$. Cột trái: $21 - 4 - 11 = 6$. Hàng dưới: $21 - 6 - 10 = 5$.
\begin{center}\lgaMagic{$4$}{$9$}{\textcolor{tmGreen!60!black}{$8$}}{$11$}{\textcolor{tmGreen!60!black}{$7$}}{\textcolor{tmGreen!60!black}{$3$}}{\textcolor{tmGreen!60!black}{$6$}}{\textcolor{tmGreen!60!black}{$5$}}{\textcolor{AnsRed}{$\mathbf{10}$}}\end{center}
\verify{Cột giữa $9 + 7 + 5 = 21$; cột phải $8 + 3 + 10 = 21$; chéo phụ $8 + 7 + 6 = 21$. Chín số $3$ đến $11$ dùng đủ, không lặp.}
\begin{warnbox}Nếu nhầm tổng ma phương là $15$ (như với các số $1$ đến $9$) thì sẽ ra $x = 4$ -- trùng với số đã có! Tổng ma phương
phụ thuộc vào các số được dùng.\end{warnbox}
Vậy chọn \textbf{(C)}.
\ans{(C) $10$}
}

% ---------------------------------------------------------------- Bài 22
\needspace{12\baselineskip}\begin{probbox}{22}{\relax}
Hai chiếc cân dưới đây đều thăng bằng. Các hình tròn nặng bằng nhau.
\begin{center}
\lgaScale{\lgaCir\,\lgaCir\,\lgaCir}{\lgaSq}\hspace{1cm}\lgaScale{\lgaSq\,\lgaCir}{\lgaWeight{800 g}}
\end{center}
Hỏi hình vuông \lgaSq\ nặng bao nhiêu gam?
\end{probbox}
\probsol{
\textbf{Phân tích đề.} Có hai đại lượng chưa biết: khối lượng một hình tròn và khối lượng hình vuông. Hai chiếc cân cho hai đẳng thức.
\begin{recallbox}Nếu nhóm vật $A$ cân bằng với nhóm vật $B$ thì ở chiếc cân khác có thể thay $A$ bằng $B$ (thay thế).\end{recallbox}
\step{1}{Gọi ẩn}
Gọi khối lượng một hình tròn là $x$ (gam), $x > 0$. Cân thứ nhất cho biết hình vuông nặng bằng ba hình tròn, nên hình vuông nặng $3x$ (gam).
\step{2}{Lập phương trình từ cân thứ hai}
Đĩa trái có hình vuông và một hình tròn: $3x + x$ (gam); đĩa phải là quả cân $800$ g. Ta có phương trình
\[ 3x + x = 800, \qquad\text{tức là}\qquad 4x = 800. \]
\step{3}{Giải phương trình}
$x = 800 : 4 = 200$. Một hình tròn nặng $200$ g.
\step{4}{Trả lời câu hỏi}
Hình vuông nặng $3x = 3\cdot 200 = 600$ (g).
\begin{warnbox}Đề hỏi khối lượng \emph{hình vuông}, không phải hình tròn. Tính ra $x = 200$ rồi dừng lại là trả lời nhầm câu hỏi.\end{warnbox}
\verify{Cân 1: $3\cdot 200 = 600$ g, bằng hình vuông \checkmark. Cân 2: $600 + 200 = 800$ g \checkmark.}
\ans{Hình vuông nặng $600$ g.}
}

% ---------------------------------------------------------------- Bài 23
\needspace{12\baselineskip}\begin{probbox}{23}{\relax}
Có bốn loại hộp: đỏ \lgaBox{AnsRed}, vàng \lgaBox{StarGold}, xanh \lgaBox{tmBlue}, tím \lgaBox{tmPurple}; các hộp cùng màu nặng bằng nhau.
Ba chiếc cân sau đều thăng bằng:
\begin{center}
\lgaScale{\lgaBox{AnsRed}\,\lgaBox{StarGold}}{\lgaBox{tmBlue}}\hspace{0.4cm}\lgaScale{\lgaBox{tmBlue}\,\lgaBox{StarGold}}{\lgaBox{tmPurple}}\hspace{0.4cm}\lgaScale{\lgaBox{AnsRed}\,\lgaBox{AnsRed}}{\lgaBox{tmBlue}}
\end{center}
Một hộp tím nặng bằng bao nhiêu hộp vàng?
\choices{$2$}{$3$}{$4$}{$5$}{$6$}
\end{probbox}
\probsol{
\textbf{Phân tích đề.} Gọi khối lượng hộp đỏ, vàng, xanh, tím là $D, V, X, T$. Ba cân cho ba đẳng thức:
\[ Đ + V = X, \qquad X + V = T, \qquad Đ + Đ = X. \]
Cần biểu diễn $T$ theo $V$.
\step{1}{So sánh cân 1 và cân 3}
Cả hai đều có đĩa phải là một hộp xanh, nên đĩa trái của chúng nặng bằng nhau: $D + V = D + D$.
Bỏ một hộp đỏ ở cả hai bên: $V = D$. \emph{Hộp đỏ nặng bằng hộp vàng.}
\step{2}{Hộp xanh}
$X = D + V = V + V = 2V$: hộp xanh nặng bằng $2$ hộp vàng.
\step{3}{Hộp tím}
$T = X + V = 2V + V = 3V$: hộp tím nặng bằng $3$ hộp vàng.
\begin{figC}
  \node at (0,0){\lgaScale{\lgaBox{tmPurple}}{\lgaBox{StarGold}\,\lgaBox{StarGold}\,\lgaBox{StarGold}}};
\end{figC}
\begin{warnbox}Đừng vội cho rằng hộp đỏ nặng hơn hộp vàng chỉ vì ``đỏ'' xuất hiện nhiều lần. Phải dùng các cân để \emph{chứng minh} quan hệ.\end{warnbox}
\verify{Chọn $V = D = 1$ kg, $X = 2$ kg, $T = 3$ kg: cân 1: $1 + 1 = 2$; cân 2: $2 + 1 = 3$; cân 3: $1 + 1 = 2$. Cả ba đều thăng bằng.}
Vậy chọn \textbf{(B)}.
\ans{(B) $3$}
}

% ---------------------------------------------------------------- Bài 24
\needspace{12\baselineskip}\begin{probbox}{24}{\relax}
\begin{minipage}[t]{0.63\linewidth}
Các số $1, 2, 3, 4, 5, 6, 7, 8$ được điền vào tám vòng tròn trên hình vuông bên cạnh (bốn ở đỉnh, bốn ở giữa cạnh), mỗi số một lần,
sao cho tổng ba số trên mỗi cạnh đều bằng $13$. Tổng bốn số ở bốn đỉnh là
\choices{$12$}{$13$}{$16$}{$20$}{$36$}
\end{minipage}\hfill
\begin{minipage}[t]{0.33\linewidth}\vspace{0pt}\centering
\lgaSqEight{}{}{}{}{}{}{}{}
\end{minipage}
\end{probbox}
\probsol{
\textbf{Phân tích đề.} Hình vuông có $4$ cạnh, mỗi cạnh gồm hai đỉnh và một số ở giữa. Cần tìm tổng bốn số ở đỉnh.
\begin{recallbox}Khi cộng tổng của tất cả các cạnh, số ở đỉnh được tính hai lần (vì thuộc hai cạnh), số ở giữa cạnh được tính một lần.\end{recallbox}
\step{1}{Cộng tổng của bốn cạnh}
Mỗi cạnh có tổng $13$, nên tổng bốn cạnh là $4\times 13 = 52$.
\step{2}{Đếm mỗi số được tính mấy lần}
Trong $52$ đó, mỗi số từ $1$ đến $8$ được tính ít nhất một lần, còn bốn số ở đỉnh được tính thêm một lần nữa:
\[ 52 = (1 + 2 + \cdots + 8) + (\text{tổng bốn đỉnh}) = 36 + (\text{tổng bốn đỉnh}). \]
\step{3}{Tính}
Tổng bốn đỉnh $= 52 - 36 = 16$.
\step{4}{Một cách điền cụ thể (để chắc chắn bài toán có lời giải)}
\begin{figC}
  \node at (0,0){\lgaSqEight{$1$}{$8$}{$4$}{$3$}{$6$}{$2$}{$5$}{$7$}};
  \node[font=\small,anchor=west,align=left] at (2.0,0){trên: $1 + 8 + 4 = 13$\\ phải: $4 + 3 + 6 = 13$\\ dưới: $6 + 2 + 5 = 13$\\ trái: $5 + 7 + 1 = 13$\\[2pt] đỉnh: $1 + 4 + 6 + 5 = 16$};
\end{figC}
\begin{warnbox}Phương án (E) $36$ là tổng của \emph{tất cả} tám số; phương án (B) $13$ là tổng một cạnh. Câu hỏi là tổng bốn số ở \emph{đỉnh}.\end{warnbox}
Vậy chọn \textbf{(C)}.
\ans{(C) $16$}
}

% ---------------------------------------------------------------- Bài 25
\needspace{12\baselineskip}\begin{probbox}{25}{\relax}
\begin{minipage}[t]{0.63\linewidth}
Điền các số $1, 2, 3, 4, 5$ vào năm ô của hình chữ thập bên cạnh, mỗi số một lần, sao cho tổng ba số trên hàng ngang bằng tổng ba số trên
cột dọc. Ô ở chính giữa có thể nhận bao nhiêu giá trị khác nhau?
\choices{$1$}{$2$}{$3$}{$4$}{$5$}
\end{minipage}\hfill
\begin{minipage}[t]{0.33\linewidth}\vspace{-6pt}\centering
\lgaCross{}{}{$?$}{}{}
\end{minipage}
\end{probbox}
\probsol{
\textbf{Phân tích đề.} Ô giữa (gọi là $c$) thuộc cả hàng ngang lẫn cột dọc. Bốn ô còn lại chia thành hai cặp: cặp trái -- phải (hàng ngang)
và cặp trên -- dưới (cột dọc).
\step{1}{Điều kiện bằng nhau}
Hàng ngang: $(\text{trái} + \text{phải}) + c$; cột dọc: $(\text{trên} + \text{dưới}) + c$. Hai tổng bằng nhau khi và chỉ khi
\[ \text{trái} + \text{phải} = \text{trên} + \text{dưới}. \]
\step{2}{Tổng của bốn số còn lại phải chẵn}
Bốn số không ở giữa có tổng $= 15 - c$ (vì $1 + 2 + 3 + 4 + 5 = 15$). Hai cặp có tổng bằng nhau nên $15 - c$ phải chia đôi được, tức là
$15 - c$ chẵn. Vì $15$ lẻ, $c$ phải \textbf{lẻ}: $c \in \{1; 3; 5\}$.
\step{3}{Kiểm tra từng giá trị lẻ}
\begin{itemize}[leftmargin=1.4em,itemsep=1pt,topsep=1pt]
\item $c = 1$: còn $2, 3, 4, 5$, mỗi cặp có tổng $14 : 2 = 7$: \ $2 + 5 = 3 + 4 = 7$. \checkmark
\item $c = 3$: còn $1, 2, 4, 5$, mỗi cặp có tổng $12 : 2 = 6$: \ $1 + 5 = 2 + 4 = 6$. \checkmark
\item $c = 5$: còn $1, 2, 3, 4$, mỗi cặp có tổng $10 : 2 = 5$: \ $1 + 4 = 2 + 3 = 5$. \checkmark
\end{itemize}
\begin{figC}
  \node at (0,0){\lgaCross{$3$}{$2$}{$1$}{$5$}{$4$}}; \node[font=\small] at (0,-1.5){$c = 1$};
  \node at (3.2,0){\lgaCross{$2$}{$1$}{$3$}{$5$}{$4$}}; \node[font=\small] at (3.2,-1.5){$c = 3$};
  \node at (6.4,0){\lgaCross{$2$}{$1$}{$5$}{$4$}{$3$}}; \node[font=\small] at (6.4,-1.5){$c = 5$};
\end{figC}
\step{4}{Vì sao $c$ chẵn không được?}
Nếu $c = 2$: bốn số còn lại $1, 3, 4, 5$ có tổng $13$ -- lẻ, không chia thành hai cặp có tổng bằng nhau. Nếu $c = 4$: còn $1, 2, 3, 5$ có tổng
$11$ -- lẻ, cũng không được.
\begin{warnbox}Đừng nghĩ ``số ở giữa phải là số đứng giữa ($3$)''. Ở đây có tới ba khả năng!\end{warnbox}
Vậy ô giữa có thể là $1$, $3$ hoặc $5$: có $3$ giá trị. Chọn \textbf{(C)}.
\ans{(C) $3$}
}

% ---------------------------------------------------------------- Bài 26
\needspace{12\baselineskip}\begin{probbox}{26}{\relax}
\begin{minipage}[t]{0.66\linewidth}
Ba con xúc xắc chuẩn (tổng hai mặt đối diện bằng $7$) được xếp chồng lên nhau trên mặt bàn như hình bên. Mặt trên cùng có $2$ chấm.
Tổng số chấm trên năm mặt bị che khuất \emph{nằm ngang} (các mặt áp vào nhau giữa hai con xúc xắc và mặt đáy áp xuống bàn) là
\choices{$14$}{$17$}{$19$}{$21$}{$23$}
\end{minipage}\hfill
\begin{minipage}[t]{0.3\linewidth}\vspace{-4pt}\centering
\begin{tikzpicture}[scale=0.62,line join=round]
  \draw[Slate,line width=1pt] (-0.4,0)--(2.0,0);
  \foreach \k/\f/\r in {0/4/1,1/5/3,2/3/1}{
    \begin{scope}[yshift=\k cm]
    \draw[Ink,line width=0.8pt,fill=white] (0,0) rectangle (1,1);
    \draw[Ink,line width=0.8pt,fill=white] (1,0)--(1.45,0.45)--(1.45,1.45)--(1,1)--cycle;
    \begin{scope}[Ink]\lgaPips{\f}\end{scope}
    \begin{scope}[Ink,cm={0.45,0.45,0,1,(1,0)}]\lgaPips{\r}\end{scope}
    \end{scope}}
  \draw[Ink,line width=0.8pt,fill=white] (0,3)--(1,3)--(1.45,3.45)--(0.45,3.45)--cycle;
  \begin{scope}[Ink,cm={1,0,0.45,0.45,(0,3)}]\lgaPips{2}\end{scope}
\end{tikzpicture}
\end{minipage}
\end{probbox}
\probsol{
\textbf{Phân tích đề.} Năm mặt bị che nằm ngang gồm: mặt đáy của con trên cùng; mặt trên và mặt đáy của con ở giữa; mặt trên và mặt đáy của
con dưới cùng. Ta không cần biết từng mặt có bao nhiêu chấm -- chỉ cần dùng ``hai mặt đối diện có tổng $7$''.
\begin{recallbox}Trên xúc xắc chuẩn, mặt trên và mặt đáy là hai mặt đối diện, nên tổng số chấm của chúng bằng $7$.\end{recallbox}
\step{1}{Con dưới cùng và con ở giữa}
Mỗi con có mặt trên và mặt đáy đều bị che, tổng mỗi cặp là $7$. Hai con: $7 + 7 = 14$.
\step{2}{Con trên cùng}
Mặt trên nhìn thấy có $2$ chấm, nên mặt đáy (đối diện) có $7 - 2 = 5$ chấm, và mặt này bị che.
\step{3}{Cộng lại}
$14 + 5 = 19$.
\begin{warnbox}(D) $21 = 3\times 7$ là quên rằng mặt trên cùng ($2$ chấm) \emph{nhìn thấy được}, không bị che. (A) $14$ là quên mặt đáy của con trên cùng
(cũng bị che, có $5$ chấm).\end{warnbox}
\verify{Sáu mặt nằm ngang của ba con có tổng $3\times 7 = 21$; trong đó chỉ mặt trên cùng ($2$ chấm) là nhìn thấy, nên năm mặt bị che có tổng $21 - 2 = 19$.}
Vậy chọn \textbf{(C)}.
\ans{(C) $19$}
}

% ---------------------------------------------------------------- Bài 27
\needspace{12\baselineskip}\begin{probbox}{27}{\relax}
\begin{minipage}[t]{0.63\linewidth}
Lưới $4\times 4$ bên cạnh cần được điền sao cho mỗi hàng, mỗi cột và mỗi ô vuông $2\times 2$ viền đậm đều có đủ các chữ số $1, 2, 3, 4$.
Chữ số ở ô $Q$ là
\choicesL{$1$}{$2$}{$3$}{$4$}{Không xác định được}
\end{minipage}\hfill
\begin{minipage}[t]{0.33\linewidth}\vspace{4pt}\centering
\lgaFour{2/1/$3$,4/1/$2$,4/2/$1$,1/3/$4$,3/4/$2$,4/4/$Q$}
\end{minipage}
\end{probbox}
\probsol{
\textbf{Phân tích đề.} Ô $Q$ nằm ở hàng $4$, cột $4$ (góc dưới phải), thuộc ô $2\times 2$ dưới phải.
\begin{recallbox}Chữ số trong một ô phải khác mọi chữ số đã có trong hàng, cột và ô $2\times 2$ chứa nó.\end{recallbox}
\step{1}{Cột $4$}
Cột $4$ đã có $2$ (hàng $1$) và $1$ (hàng $2$). Hai ô còn lại -- hàng $3$ cột $4$ và ô $Q$ -- phải là $3$ và $4$ theo một thứ tự nào đó.
\step{2}{Hàng $3$ quyết định}
Hàng $3$ đã có chữ số $4$ (ở cột $1$), nên hàng $3$ cột $4$ \emph{không} thể là $4$. Vậy hàng $3$ cột $4$ là $3$, và $Q = 4$.
\step{3}{Kiểm tra với ô $2\times 2$ dưới phải và hàng $4$}
Ô $2\times 2$ dưới phải chứa $3$ (hàng $3$ cột $4$), $2$ (hàng $4$ cột $3$), $Q = 4$ và ô còn lại (hàng $3$ cột $3$) là $1$ -- đủ $1, 2, 3, 4$.
Hàng $4$ có $2$ và $4$, không lặp.
\step{4}{Hoàn thành lưới (để chắc chắn không mâu thuẫn)}
\begin{figC}
  \node at (0,0){\lgaFour{1/1/$1$,2/1/$3$,3/1/$4$,4/1/$2$,1/2/$2$,2/2/$4$,3/2/$3$,4/2/$1$,1/3/$4$,2/3/$2$,3/3/$1$,4/3/$3$,1/4/$3$,2/4/$1$,3/4/$2$,4/4/\textcolor{AnsRed}{$\mathbf{4}$}}};
  \node[font=\small,anchor=west,align=left] at (1.8,0){Hàng $1$: $1\,3\,4\,2$\\ Hàng $2$: $2\,4\,3\,1$\\ Hàng $3$: $4\,2\,1\,3$\\ Hàng $4$: $3\,1\,2\,4$};
\end{figC}
\verify{Mọi hàng, cột và ô $2\times 2$ đều có đủ $1, 2, 3, 4$; máy tính kiểm tra thấy lưới chỉ có duy nhất cách điền này.}
\begin{warnbox}Chỉ nhìn cột $4$ thì thấy $Q$ có thể là $3$ hoặc $4$ và dễ chọn (E). Hãy nhìn thêm hàng $3$ -- một vùng khác chứa ô ``hàng xóm'' của $Q$.\end{warnbox}
Vậy chọn \textbf{(D)}.
\ans{(D) $4$}
}

% ---------------------------------------------------------------- Bài 28
\needspace{12\baselineskip}\begin{probbox}{28}{\relax}
Ở lớp 7A, \emph{tất cả} các bạn trong câu lạc bộ cờ vua đều biết bơi. Bạn Leo không biết bơi. Bạn Nora biết bơi. Khẳng định nào sau đây
\textbf{chắc chắn} đúng?
\choicesL{Leo ở trong câu lạc bộ cờ vua.}{Leo không ở trong câu lạc bộ cờ vua.}{Nora ở trong câu lạc bộ cờ vua.}{Nora không ở trong câu lạc bộ cờ vua.}{Mọi bạn biết bơi trong lớp đều ở câu lạc bộ cờ vua.}
\end{probbox}
\probsol{
\textbf{Phân tích đề.} Câu ``tất cả các bạn trong CLB cờ vua đều biết bơi'' có dạng ``mọi A đều là B'' (A: ở CLB cờ vua, B: biết bơi).
\begin{recallbox}Từ ``mọi A đều là B'' suy ra được ``không là B thì không là A''; nhưng \emph{không} suy ra được ``là B thì là A''.\end{recallbox}
\step{1}{Vẽ sơ đồ}
\begin{figC}
  \draw[tmBlue,line width=1pt,fill=tmBlue!8] (0.9,1.1) ellipse (1.9 and 1.2);
  \draw[tmOrange,line width=1pt,fill=tmOrange!15] (0.3,1.0) ellipse (0.95 and 0.6);
  \node[font=\scriptsize,text=tmOrange!70!black] at (0.3,1.25){CLB cờ vua};
  \node[font=\scriptsize,text=tmBlue!80!black] at (1.95,1.6){biết bơi};
  \fill[tmGreen!60!black] (1.95,0.95) circle (2.2pt) node[below,font=\scriptsize]{Nora?};
  \fill[tmGreen!60!black] (0.3,0.85) circle (2.2pt) node[below,font=\scriptsize]{Nora?};
  \fill[AnsRed] (3.6,1.1) circle (2.2pt) node[right,font=\scriptsize]{Leo};
\end{figC}
Vòng ``CLB cờ vua'' nằm trọn trong vòng ``biết bơi''.
\step{2}{Leo}
Leo không biết bơi, nên Leo nằm \emph{ngoài} vòng lớn, do đó chắc chắn nằm ngoài vòng nhỏ: \textbf{Leo không ở CLB cờ vua}. Nếu Leo ở CLB thì
Leo phải biết bơi -- trái với đề (phản chứng).
\step{3}{Nora}
Nora biết bơi: Nora nằm trong vòng lớn, nhưng có thể ở trong hoặc ngoài vòng nhỏ. Vậy (C) và (D) đều chỉ \emph{có thể} đúng, không chắc chắn.
\step{4}{Câu (E)}
(E) đảo ngược câu đề bài. Đề không nói gì về những bạn biết bơi nhưng không chơi cờ, nên (E) không chắc chắn đúng.
\begin{warnbox}Lỗi rất phổ biến: từ ``người chơi cờ đều biết bơi'' suy ra ``người biết bơi đều chơi cờ''. Hãy nghĩ một ví dụ: ``mọi con mèo đều có
bốn chân'' không có nghĩa ``mọi con vật bốn chân đều là mèo''!\end{warnbox}
Vậy chọn \textbf{(B)}.
\ans{(B) Leo không ở trong câu lạc bộ cờ vua}
}

% ---------------------------------------------------------------- Bài 29
\needspace{12\baselineskip}\begin{probbox}{29}{\relax}
Chú Cáo nói dối vào các ngày thứ Hai, thứ Ba, thứ Tư và nói thật vào các ngày còn lại trong tuần. Một hôm, Cáo nói: ``Hôm qua tôi đã nói dối.''
Hôm đó có thể là ngày nào?
\choicesL{Chỉ có thể là thứ Hai.}{Chỉ có thể là thứ Năm.}{Thứ Hai hoặc thứ Năm.}{Thứ Tư hoặc thứ Năm.}{Chủ nhật hoặc thứ Hai.}
\end{probbox}
\probsol{
\textbf{Phân tích đề.} Cáo nói được câu ``Hôm qua tôi đã nói dối'' khi:
\begin{itemize}[leftmargin=1.4em,itemsep=0pt,topsep=1pt]
\item hôm nay nói thật và câu đúng: hôm nay thật, hôm qua dối; hoặc
\item hôm nay nói dối và câu sai: hôm nay dối, hôm qua thật.
\end{itemize}
Tức là: \textbf{hôm qua và hôm nay khác loại}.
\begin{recallbox}Hôm qua của thứ Hai là Chủ nhật (tuần quay vòng).\end{recallbox}
\step{1}{Lập bảng}
\begin{center}
\setlength{\tabcolsep}{4pt}
\begin{tabular}{|l|c|c|c|c|c|c|c|}
\hline
\rowcolor{tmBlue!10} Hôm nay & T2 & T3 & T4 & T5 & T6 & T7 & CN \\ \hline
Hôm nay Cáo & D & D & D & T & T & T & T \\ \hline
Hôm qua Cáo & T (CN) & D & D & D & T & T & T \\ \hline
Khác loại? & \textbf{có} & -- & -- & \textbf{có} & -- & -- & -- \\ \hline
\end{tabular}
\end{center}
\step{2}{Kiểm tra hai ngày tìm được}
\emph{Thứ Hai:} Cáo nói dối; hôm qua (Chủ nhật) Cáo nói thật, nên câu ``hôm qua tôi nói dối'' sai -- đúng là một lời nói dối. \checkmark\\
\emph{Thứ Năm:} Cáo nói thật; hôm qua (thứ Tư) Cáo nói dối, nên câu đó đúng -- đúng là lời nói thật. \checkmark
\begin{warnbox}Thứ Tư: Cáo nói dối, hôm qua (thứ Ba) cũng nói dối, nên câu ``hôm qua tôi đã nói dối'' là câu \emph{đúng}. Người đang nói dối không
nói câu đúng, nên thứ Tư bị loại -- đây là bẫy của phương án (D).\end{warnbox}
Vậy chọn \textbf{(C)}.
\ans{(C) Thứ Hai hoặc thứ Năm}
}

% ---------------------------------------------------------------- Bài 30
\needspace{12\baselineskip}\begin{probbox}{30}{\relax}
Bốn bạn Amy, Ben, Chloe, David thi chạy, không có hai bạn nào về cùng lúc. Ba bạn đứng xem dự đoán:
\begin{itemize}[leftmargin=1.6em,itemsep=0pt,topsep=2pt]
\item Henry: ``Amy về nhất, Ben về nhì.''
\item Kate: ``Chloe về nhất, David về thứ tư.''
\item Lucy: ``Ben về thứ ba, David về nhì.''
\end{itemize}
Kết quả là mỗi bạn dự đoán đúng \emph{đúng một} ý. Hãy tìm thứ tự về đích của bốn bạn.
\end{probbox}
\probsol{
\textbf{Phân tích đề.} Mỗi dự đoán có hai ý, đúng một ý đúng. Ta xét trường hợp theo dự đoán của Henry.
\begin{recallbox}Câu ``đúng một nửa'': ý này đúng thì ý kia sai, ý này sai thì ý kia đúng.\end{recallbox}
\step{1}{Trường hợp 1: ``Ben về nhì'' đúng (và ``Amy về nhất'' sai)}
Xét dự đoán của Lucy. Ben về nhì nên ``Ben về thứ ba'' sai; vậy ý còn lại ``David về nhì'' phải đúng. Nhưng Ben đã về nhì, David không thể
cũng về nhì. \textbf{Mâu thuẫn} -- loại trường hợp này.
\step{2}{Trường hợp 2: ``Amy về nhất'' đúng (và ``Ben về nhì'' sai)}
\begin{itemize}[leftmargin=1.4em,itemsep=1pt,topsep=1pt]
\item Kate: Amy về nhất nên ``Chloe về nhất'' sai $\Rightarrow$ ``David về thứ tư'' đúng.
\item Lucy: David về thứ tư nên ``David về nhì'' sai $\Rightarrow$ ``Ben về thứ ba'' đúng.
\item Còn lại vị trí thứ nhì cho Chloe.
\end{itemize}
Thứ tự: Amy (nhất), Chloe (nhì), Ben (ba), David (tư).
\step{3}{Thử lại}
\begin{center}
\begin{tabular}{|l|c|c|}
\hline
\rowcolor{tmBlue!10} Người dự đoán & Ý thứ nhất & Ý thứ hai \\ \hline
Henry & Amy nhất: \textcolor{tmGreen!70!black}{đúng} & Ben nhì: \textcolor{AnsRed}{sai} \\ \hline
Kate & Chloe nhất: \textcolor{AnsRed}{sai} & David tư: \textcolor{tmGreen!70!black}{đúng} \\ \hline
Lucy & Ben ba: \textcolor{tmGreen!70!black}{đúng} & David nhì: \textcolor{AnsRed}{sai} \\ \hline
\end{tabular}
\end{center}
\begin{tipbox}Chọn dự đoán để xét trường hợp sao cho mỗi trường hợp ``kéo theo'' được nhiều ý khác. Ở đây ý ``Ben nhì'' va chạm ngay với
``David nhì'' của Lucy, nên trường hợp 1 bị loại rất nhanh.\end{tipbox}
\begin{parentbox}Có thể cùng con dùng bốn mẩu giấy ghi tên, xếp thành hàng ``nhất -- nhì -- ba -- tư'' và đổi chỗ để thử. Sau đó hỏi con:
``Mình có thử hết các trường hợp chưa? Làm sao chắc chắn không còn đáp án khác?'' -- câu trả lời chính là Bước 1.\end{parentbox}
\ans{Amy nhất, Chloe nhì, Ben ba, David tư.}
}
